\subsection*{幂等等式观察哪些状态}

用 \(s\) 表示系统状态，用 \(f\) 表示执行一个操作所定义的状态变化。数学上的幂等要求，对约定状态集合内的每个 \(s\)，都有
\[
 f(f(s))=f(s).
\]
内层 \(f(s)\) 是执行一次的状态，外层 \(f\) 是紧接着再执行同一变换；等号要求执行两次与执行一次的最终状态相同。状态包含哪些字段，决定这句话实际约束什么。

只看余额 \(b\)，把余额固定为 7000 的变换记为 \(a(b)=7000\)，就有 \(a(a(b))=7000=a(b)\)。这个等式没有安排两次 \(a\) 之间插入别的操作。让 \(g(b)=b-1000\) 表示另一项扣款，先 \(a\) 后 \(g\) 再 \(a\)，从 10000 演算得到 7000、6000、7000；原本希望原 \(a\) 完成后只执行 \(g\)，结果应为 6000。

若把订单和流水也放进 \(s\)，每次 \(a\) 都新增一条流水，则第二次 \(a\) 还会增加记录。即使余额保持 7000，完整状态已经与一次执行不同。讨论幂等需要先列出观察范围，才能判断等号是否确实成立。

\begin{figure}[htbp]
\centering
\begin{tikzpicture}[x=1cm,y=1cm]
\node[booknote,text width=13cm] at (7,.45) {$a(b)=7000$；$g(b)=b-1000$；共同初始余额为 10000};
\node[booknode,text width=4.4cm] (adj) at (3,-.7) {连续重复\\$10000\stackrel{a}{\longrightarrow}7000\stackrel{a}{\longrightarrow}7000$};
\node[booknode,text width=4.4cm] (eq) at (11,-.7) {$a(a(10000))=a(10000)$\\余额投影上相同};
\draw[bookarrow] (adj)--(eq);
\node[booknode,text width=4.4cm] (inter) at (3,-2.4) {中间插入 $g$ 后重复 $a$\\$7000\stackrel{g}{\longrightarrow}6000\stackrel{a}{\longrightarrow}7000$};
\node[bookresult,text width=4.4cm] (diff) at (11,-2.4) {原 $a$ 完成后只执行 $g$\\最终应为 6000};
\draw[bookalert] (inter)--(diff);
\node[booknote,text width=13cm] at (7,-3.8) {连续重复等式没有检查 $a\to g\to a$ 的交错；迟到请求需要恢复原结果};
\end{tikzpicture}
\caption{幂等等式比较紧邻重复，插入其他操作需要另查执行历史。}
\label{fig:ch04-adjacent-idempotence}
\end{figure}


可以给观察范围一个明确记号。设完整业务状态为 \(s=(b,P,O,D,M)\)：\(b\) 是余额，\(P\)、\(O\)、\(D\) 分别是购买、订单与扣款记录集合，\(M\) 是身份绑定与原结果。只取余额的观察记为 \(\pi_b(s)=b\)。这种从完整状态中选取一部分信息的映射称为投影，它用来说明当前比较究竟保留了什么、忽略了什么。

假设变换 \(F\) 每次把余额设为 7000，并新添一条 3000 分流水。从初始 \(s_0\) 执行一次得到 \(s_1\)，再执行一次得到 \(s_2\)。于是
\[
 \pi_b(s_1)=\pi_b(s_2)=7000,
 \qquad |D(s_1)|=1,\quad |D(s_2)|=2.
\]
这里 \(|D(s)|\) 表示状态 \(s\) 中扣款记录的条数。余额投影相等，完整状态仍有 \(s_1\ne s_2\)。因为投影舍弃了流水，投影上的相等无法推出完整状态相等。这个反例说明只测试余额没有继续下降，可能遗漏第二条扣款记录。

\begin{figure}[htbp]
\centering
\begin{tikzpicture}[x=1cm,y=1cm]
\node[booknode,text width=4.8cm] (s1) at (3,0) {执行一次后的完整状态 $s_1$\\余额 7000；扣款 1 条};
\node[bookresult,text width=4.8cm] (s2) at (11,0) {执行两次后的完整状态 $s_2$\\余额 7000；扣款 2 条};
\node[booknode,text width=4.8cm] (b1) at (3,-2) {只取余额：$\pi_b(s_1)=7000$};
\node[booknode,text width=4.8cm] (b2) at (11,-2) {只取余额：$\pi_b(s_2)=7000$};
\draw[bookarrow] (s1)--node[left,font=\footnotesize]{投影}(b1);
\draw[bookarrow] (s2)--node[right,font=\footnotesize]{投影}(b2);
\draw[bookalert] (s1)--node[above,font=\footnotesize]{$s_1\ne s_2$}(s2);
\draw[bookarrow] (b1)--node[above,font=\footnotesize]{相等}(b2);
\node[booknote,text width=13cm] at (7,-3.4) {投影舍弃流水条数；余额相等只说明这一项观察相同};
\end{tikzpicture}
\caption{观察余额会隐藏新增流水，投影上的相等不足以推出完整状态相等。}
\label{fig:ch04-state-projection}
\end{figure}


还要检查投影是否足以预测后续处理。两个状态即使都有余额 7000，其中一个保留 S 的成功结果，另一个没有身份记录，重试 S 时会走不同分支。仅给定 \(b=7000\)，无法决定应当返回 P1 还是尝试新购买。要把整个服务简化成余额上的函数，需要证明被省略的信息不会影响后续结果；本例的 S 与原结果恰好影响恢复，因而应保留在模型状态中。

本章的恢复操作可以记为 \(f_S\)，输入带着固定 S 和已绑定输入。状态中若没有 S，它尝试完成购买并保存结果；若已有同参终态，它读取结果，不再改变购买效果。对同一完整、有效且保留中的 S，已经提交 \(f_S\) 后再执行 \(f_S\)，账户、购买、订单和资金记录都保持原样。每次请求仍可新增访问日志；本章把通信诊断日志排除在购买效果范围之外，并保留对每次尝试的独立记录。

严格地说，服务还会返回原业务结果 \(R(S)\)。可以把一次已经确定的处理写成 \(F_S(s)=(s',R(S))\)，其中 \(s'\) 是提交后的业务状态。同参恢复时要求
\[
 F_S(s')=(s',R(S)).
\]
这个表达同时检查状态保持与结果恢复。若只检查状态没有再次扣款，却在恢复时重新用当前余额生成拒绝，就可能保住资金效果而返回错误结论。原结果需要由 \(M\) 一起保存，才能使第二项也成立。

这些表达依赖 S 的输入绑定。相同 S 携带不同参数，服务走冲突分支，不在上述同参恢复公式的前提之内。它们也依赖身份记录在允许重试时保留、已提交事实未丢失。公式描述已经确定的逻辑操作及其恢复，锁等待、超时和明确回滚仍由执行协议处理；一次尚未确定结果的尝试不能直接当作完成的 \(F_S\)。

\subsection*{把迟到赋值放回交错历史}

第四节的 S1 先扣 3000，S2 再扣 1000，流水合计要求余额为 \(10000-3000-1000=6000\)。若 S1 的恢复重新执行固定赋值 7000，D2 仍在，账户却多出 1000，资金核对关系被破坏。删除 D2 来凑平余额，又会删除一项真实成交；正确恢复应读取 S1 的完成结果，保留 S2 已经造成的变化。

\begin{figure}[htbp]
\centering
\begin{tikzpicture}[x=1cm,y=1cm]
\node[booknode,text width=3.7cm] (a) at (2,0) {S1 写固定值 7000\\10000 → 7000};
\node[booknode,text width=3.7cm] (g) at (7,0) {S2 合法扣 1000\\7000 → 6000};
\node[bookresult,text width=3.7cm] (aa) at (12,0) {S1 重试再写 7000\\6000 → 7000};
\draw[bookarrow] (a)--(g);\draw[bookalert] (g)--(aa);
\node[booknode,text width=3.7cm] (d1) at (2,-1.9) {D1：扣 3000};
\node[booknode,text width=3.7cm] (d2) at (7,-1.9) {D2：扣 1000};
\node[bookresult,text width=3.7cm] (bad) at (12,-1.9) {D2 仍在\\余额中的变化被覆盖};
\draw[bookarrow] (a)--(d1);\draw[bookarrow] (g)--(d2);\draw[bookalert] (aa)--(bad);
\node[booknote,text width=13cm] at (7,-3.2) {流水核对要求 $10000-3000-1000=6000$；实际余额 7000，相差 1000};
\end{tikzpicture}
\caption{连续重复的固定赋值，在另一项交易插入后会覆盖已提交的资金变化。}
\label{fig:ch04-assignment-overwrite}
\end{figure}


RIFL 用共享数值给出了更小的反例：初始为 1，客户端 B 写入 2 后服务崩溃，B 没收到响应；客户端 A 读到 2，再写入 3；B 重试旧写入，A 最后又读到 2\cite{rifl}。每次赋值紧邻重复都幂等，整段历史却无法把 B 的同一逻辑写入只安排一次。

这里所说的线性化，要求每项逻辑操作像是在调用开始与返回之间的某个时刻一次生效，并尊重已经返回以后才开始的实际先后；排出的顺序还要解释每次读取。A 第一次读到 2，要求 B 的写入在它以前；A 随后写 3，最后又读到 2，则需要写 2 发生在写 3 以后。初始为 1、期间只有这些写入时，B 在顺序中只出现一次，无法同时解释这两个观察值。

保存 B 的完成结果、重试时读取原结果，就能保留 A 已经写入的 3。这个例子同时核对中间已经返回的读取与最后状态；只比较最终一个数值，会漏掉前后观察对同一逻辑操作提出的矛盾要求。

\subsection*{交换律检查操作次序}

两个变换 \(f\) 与 \(g\) 具有交换性，是指对于允许的状态 \(s\)，
\[
 f(g(s))=g(f(s)).
\]
右边表示先 \(f\) 后 \(g\)，左边表示先 \(g\) 后 \(f\)，等式比较调换次序后的结果。它与连续执行同一变换的幂等性是两项分别可以检验的性质。

暂时只看整数余额，不加余额不足判断，扣 3000 记为 \(d(b)=b-3000\)，扣 1000 记为 \(g(b)=b-1000\)。两种次序都从 10000 得到 6000，因此在这个简化数值模型中可交换；连续执行 \(d\) 两次得到 4000，它又不满足幂等。固定赋值 \(a\) 自身幂等，但 \(a(g(10000))=7000\)，\(g(a(10000))=6000\)，它与该扣款不交换。

\begin{figure}[htbp]
\centering
\begin{tikzpicture}[x=1cm,y=1cm]
\node[booknote,text width=13cm] at (7,.45) {只看整数余额：$a(b)=7000$，$d(b)=b-3000$，$g(b)=b-1000$};
\node[booknode,text width=4.8cm] (ag) at (3,-.7) {先 $g$ 后 $a$\\$a(g(10000))=7000$};
\node[bookresult,text width=4.8cm] (ga) at (11,-.7) {先 $a$ 后 $g$\\$g(a(10000))=6000$};
\draw[bookalert] (ag)--node[above,font=\footnotesize]{不同}(ga);
\node[booknode,text width=4.8cm] (dg) at (3,-2.5) {先 $d$ 后 $g$\\$g(d(10000))=6000$};
\node[booknode,text width=4.8cm] (gd) at (11,-2.5) {先 $g$ 后 $d$\\$d(g(10000))=6000$};
\draw[bookarrow] (dg)--node[above,font=\footnotesize]{相同}(gd);
\node[booknote,text width=13cm] at (7,-4) {加入余额条件、订单和当次余额后，需要在扩展后的状态范围重新比较};
\end{tikzpicture}
\caption{固定赋值可以幂等而不与扣款交换，纯减法可以交换而不幂等。}
\label{fig:ch04-commutativity}
\end{figure}


把业务规则加回来，交换性可能改变。设初始余额为 3500，S1 申请扣 3000，S2 申请扣 1000，都要求余额足够。先执行 S1，S1 成功，余 500，S2 以余额不足结束；先执行 S2，S2 成功，余 2500，S1 以余额不足结束。最终余额和成功订单都不同。两个次序可以分别满足授权与余额条件，但它们没有产生同一状态。

即使初始 10000 足够两项都成功，把当次扣后余额纳入记录后，两种次序的完整状态也可能不同。先 S1 后 S2，D1 当次余额 7000，D2 为 6000；先 S2 后 S1，D2 当次余额 9000，D1 为 6000。最终余额都为 6000，流水里的历史时点仍有区别。因此声称两个操作交换时，要明确只比较最终余额，还是比较完整订单、流水与原结果。

交换等式只比较执行次序产生的状态，它没有检查发起人是否有权扣款。两笔未经授权的纯减法仍可在整数上交换，却都违反购买权限；允许负数的减法也能满足交换等式，却可能使余额跌到零以下。购买入口需要先核对账户权限、操作许可和报价金额，扣款时检查可用余额，并把购买、订单与流水按同一金额关联保存。这些步骤各自检查一项业务事实，交换性只用于判断已经允许的操作调换次序以后，所选观察范围是否保持相同。

如果某类业务操作在完整观察范围内确实交换，就能容许更多重排。例如仅统计两个不同、已验证事件各自增加的计数，最终和可能与到达次序无关。但重复事件依然可能多加一次，事件是否经过授权也仍然需要检查。每个事件送达几次、是否按顺序送达，需要分别写入处理规则。

当前购买还要向调用者返回订单号和当次余额，因此会留下次序证据。系统可以让数据库在并发竞争中确定实际成立的记录，再保存各笔结果。客户端无需在发送之前知道两项购买谁先取得账户行；它在恢复以后应取得与实际记录一致的原结果。不同合法调度产生不同、但各自符合规则的结果，与同一已提交购买在重试后被改写，需要分别判断。

\subsection*{沿记录关系逐项检查正确性}

不变量是系统在每个允许的已提交状态里都应保持的条件。这里可以要求余额非负，要求每笔成功购买恰好有一张对应订单和一条对应扣款，要求流水金额与成交快照相符，还要求同一 S 最多形成一笔成功购买。每项条件都要对应可检查的记录关系和执行步骤。

余额非负由条件 UPDATE 配合表级检查保护；它本身允许初始 10000 被同一意图扣两次，最后剩 4000。同一 S 的唯一约束限制身份记录数量；它本身不要求订单金额正确，也不阻止取得身份后的代码错误地扣 4000。共同事务提交使所列修改一起成立；它同样可以一起提交一组内容错误的记录。应用还要验证 Q7、使用 3000、保存正确关联，并限制所有写入路径。

在只允许本地购买扣款的范围内，可以把资金核对关系写得更具体。令初始余额为 \(b_0\)，\(D\) 为已提交扣款集合，\(a(d)>0\) 是流水 \(d\) 的金额，则每个已提交状态需要满足
\[
 b=b_0-\sum_{d\in D}a(d),\qquad b\geq 0.
\]
若新增合法充值，还要把已提交充值金额加到等式右边；当前模型的动作集合随之改变。先明确允许哪些变动，才能根据流水核对账户。以本节的变换 $F$ 连续执行两次为例，初始余额为 10000，两条流水合计 6000，资金核对式要求余额为 4000，实际余额却仍为 7000，因而违反这项关系。

不错误合并需要检查另一个方向。S1 与 S2 是 Alice 的两项明确购买，即使输入相同也分别建立记录；用商品内容去重，会使这一要求失败。同 S 异参被拒绝，既保护原购买的解释，也让客户端发现身份误复用。最多一次效果、两项购买各自存在、原结果稳定和账户资金关系，需要同时在同一条执行历史中核对。

\begin{figure}[htbp]
\centering
\begin{tikzpicture}[x=1cm,y=1cm]
\node[booknode,text width=4.8cm] (id) at (3,0) {同 S 至多一笔购买\\完整身份唯一、记录保留};
\node[booknode,text width=4.8cm] (input) at (11,0) {两次合法购买分别受理\\输入绑定按身份核对};
\node[booknode,text width=4.8cm] (money) at (3,-1.9) {余额非负与资金核对\\条件更新、正金额、共同提交};
\node[booknode,text width=4.8cm] (result) at (11,-1.9) {恢复原终态与原编号\\保存并读取原结果};
\node[bookresult,text width=10.5cm] (all) at (7,-3.8) {检查同一条已提交历史\\S → P → O、D；金额、余额与原结果共同一致};
\draw[bookarrow] (id.south)--(money.north);
\draw[bookarrow] (input.south)--(result.north);
\draw[bookarrow] (money.south)--(all.north west);
\draw[bookarrow] (result.south)--(all.north east);
\end{tikzpicture}
\caption{身份、输入、资金和结果各有检查对象，业务正确性需要这些关系共同成立。}
\label{fig:ch04-invariant-checks}
\end{figure}


\begin{table}[htbp]
\centering\small
\begin{tabularx}{\linewidth}{@{}p{2cm}XX@{}}
\toprule
检查对象 & 本例需要成立的条件 & 对应机制\\
\midrule
同一 S & 至多一笔已提交购买关联 & 完整非空组合唯一约束\\
新 S 与旧 S & 相同商品也能分别受理 & 身份来自确认动作\\
同 S 的输入 & 已保存绑定保持不变 & 按保存版本规范化并核对\\
可用余额 & 每笔实际扣款时金额足够 & 条件 UPDATE 与非负检查\\
订单与扣款 & 与购买快照共同成立 & 同一事务及业务关联检查\\
恢复结果 & 回到原终态与原编号 & 身份关联持久原结果\\
\bottomrule
\end{tabularx}
\caption{各项业务要求分别落实到可以检查的记录关系与步骤。}
\end{table}

还可以检查同一 S 的拒绝结果是否也只保留一个。余额不足提交以后，后来的同参请求应读回这条拒绝；如果另一个路径把它改为成功，同一购买的含义就变化了。终态稳定是本章所选接口规则的一部分。退款、重新购买和人工纠正应以新的业务动作表达，并引用原记录，保留已经返回或被查询过的事实。

可以按两个同 S 请求的起点再做一遍完整手工检查。开始时 S 无记录，A、B 都有同一已确认输入。若 A 先插入，B 在唯一检查等待。A 提交成功，B 读取原成功；A 提交余额不足，B 读取原拒绝；A 回滚，B 才有机会插入并继续。A 回滚后又恢复，它同样重新尝试插入，可能等待 B，再读取 B 最后提交的终态。分支取决于数据库实际提交的记录，候选编号的产生次序或客户端的响应次序只描述尝试经过。

若 B 在等待中超时并回滚，A 继续完成，后来请求 C 仍使用原 S，读取 A 的结果。若 A 与 B 都明确回滚，后来的 C 可以第一次完成这项操作。若某个连接失效而提交是否发生尚不明确，C 依旧靠唯一检查与记录读取分辨，遇到尚未结束的竞争就接受等待或暂时错误。这些分支都以同一 S 为共同对象，恢复代码按原身份继续，无需预先选定哪一个候选购买号应当成功。

\subsection*{并发控制还留下哪些故障位置}

把完整身份保存在购买表中，还要关注哪些字段允许更新。订单状态可以随履约推进发生变化，但用于识别 S 的四列与原绑定输入在终态后保持稳定。若管理员把 K1 改成 K9，原 K1 恢复就会查不到；若把原输入改成当前购物车，后来的冲突判断也会失去依据。可以通过写入权限、受控更新语句和审计限制这些修改，让身份约束保护的对象在整个恢复期内保持同一含义。

查询时也要把四列完整写出。只按 K1 读取可能拿到另一个账户的购买，只按账户与 K1 读取可能混入其他动作或会话。候选插入使用完整 S，恢复查询也使用相同完整 S，两者才能指向同一组唯一记录。输入比较从找到的记录读取原规范化版本；服务当前默认版本只用于尚未绑定的新操作。先找到旧身份，再按旧版本解释输入，使升级后的恢复仍沿用原绑定。

所有分支都可留下本次尝试的诊断记录，例如 R2 等待了多久、候选 P2 是否插入、最后返回了哪条原购买。日志中的候选编号应明确标记其尚未成立，避免排查者把两次候选分配误数成两次成交。核对效果时，从已提交的购买、订单和资金流水取证；日志帮助重建执行经过，持久业务记录说明哪些变化最终成立。

真实服务还会在每条语句之间退出、在提交前后断开连接，或在结果读取之后丢失响应。这些中断需要接回前面的等待、提交与回滚分支，再检查全部关联记录；同一身份的数量上限还要与原结果恢复一起成立。

就当前规则而言，可以先沿两种已提交事实演算。S 存在同参终态时，恢复读取保存的 P1 及结果；S 尚无已提交终态时，恢复仍须经过唯一插入，并接受可能的等待、回滚或原事务提交。错误出口保留同一身份的恢复路径，新请求取得 S 后也把原结果与资金变化放进共同事务，使另一连接能够读取完整事实。

RIFL 把调用身份、完成结果、持久保存和与操作变化的原子关系共同纳入重复调用处理机制\cite{rifl}。本章将它们缩小到一个数据库事务域，并具体展示并发者怎样在同一唯一键前等待。副本丢失提交、外部支付或身份过期清理，会改变原记录能否找到、谁仍可提交，需要按各自故障条件重新核对。

Alice 现在可以有 P1 成功以后响应丢失、S2 随后合法消费、R2 最后恢复 P1 的历史：当前余额保持 6000，P1 的原结果保留当次 7000，S2 的订单与扣款仍在。同一身份把迟到请求接回原购买，结果记录则要继续保留足够的成交内容，使服务在后来的余额、价格与接口表示变化以后仍能准确回答 P1。
\FloatBarrier
